\section{Patient Rights, Costs, and H. R. 9510 v5}

\subsection{The premise, stated as legislation would state it}

The cancer patient, not the doctor, not the nurse, not the trial sponsor, not
the Institutional Review Board, and not the regulator, is the priority
participant in a United States oncology clinical trial. That premise is the
subject of the author's proposed legislation, cited throughout this paper as
H. R. 9510 v5 \cite{hr9510billv5}, which supersedes the earlier H. R. 9501 to
H. R. 9507 numbering used in the antecedent analysis \cite{paibillprior}. The
case for enactment is developed separately \cite{congress9510, fedlawhr9510}.

The bill matters to this paper for a specific reason. Six of its provisions
describe things this protocol already does voluntarily. The difference between
doing them voluntarily and doing them under statute is that a voluntary
commitment can be withdrawn by the party that made it, and a statutory one
cannot. A participant reading the table below should understand the right-hand
column as a description of current practice at one site, and the left-hand
column as what would make it general.

\vspace{0.30em}

\begin{xltabular}{\textwidth}{L{4.3cm} Y}
\toprule
\textbf{What the bill would establish} & \textbf{What this protocol already does voluntarily} \\
\midrule
\endfirsthead
\multicolumn{2}{@{}l}{\footnotesize\itshape Table continued from the previous page.}\\
\toprule
\textbf{What the bill would establish} & \textbf{What this protocol already does voluntarily} \\
\midrule
\endhead
\midrule
\multicolumn{2}{r@{}}{\footnotesize\itshape Table continued on the next page.}\\
\endfoot
\bottomrule
\endlastfoot
Participant self-selection of trials, with round-the-clock booking & Direct booking channel with a 72-hour provisional hold and a five-business-day target, \S 6.4 and Figure 15 \\
Participant choice of robot and of autonomy level & Autonomy level set per operative step and recorded in the consent, \S 6.5 and Figure 16 \\
Participant procedural modification authority & Any excluded step is performed manually, and the exclusion is logged and reported with the outcome \\
Verification before generation, with published gates & Phase 0 gate at 1000 simulations across two frameworks, USL at least 7.0, \S 5.5 and Figure 13 \\
Real-time participant to sponsor communication & A named human at the sponsor answers within one business day, logged verbatim in your record \\
Participant health data self-custody & Per-store read lists published, and your own copy available at any time, \S 9.3 and Figure 25 \\
\end{xltabular}

\vspace{0.30em}

\subsection{Your rights during the study, enumerated}

Some of the rights below are already required by federal regulation. Others go
beyond the regulatory minimum. The distinction is marked, because a document
that presented a legal minimum as a favour would deserve exactly the distrust
the surveyed literature reports.

\vspace{0.30em}

\begin{xltabular}{\textwidth}{L{5.4cm} C{2.3cm} Y}
\toprule
\textbf{Right} & \textbf{Source} & \textbf{Note} \\
\midrule
\endfirsthead
\multicolumn{3}{@{}l}{\footnotesize\itshape Table continued from the previous page.}\\
\toprule
\textbf{Right} & \textbf{Source} & \textbf{Note} \\
\midrule
\endhead
\midrule
\multicolumn{3}{r@{}}{\footnotesize\itshape Table continued on the next page.}\\
\endfoot
\bottomrule
\endlastfoot
To be informed of the study's purpose, procedures, risks, and alternatives before consenting & 21 CFR \S 50.25 & The regulatory minimum \cite{ecfr2026part50, NCIConsent24} \\
To withdraw at any time without penalty & 21 CFR \S 50.25 & The regulatory minimum \\
To be told of significant new findings affecting your willingness to continue & 21 CFR \S 50.25 & The regulatory minimum \\
To contact the IRB directly about the conduct of the study & 21 CFR \S 50.25 & The regulatory minimum \\
To be shown the Phase 0 simulation sign-off at the baseline visit & This protocol & Beyond the minimum \\
To set the autonomy level accepted for each operative step & This protocol & Beyond the minimum \\
To be re-consented before any software version change takes effect & This protocol & Beyond the minimum \\
To request your own audit-trail extract, with a verification tool & This protocol & Beyond the minimum \\
To receive your own day-30 and day-90 results directly, not on request & This protocol & Beyond the minimum \\
To be told the operating surgeon's case count on this platform & This protocol & Beyond the minimum \\
To have your treating oncologist notified on every branch of every route & This protocol & Beyond the minimum \\
\end{xltabular}

\vspace{0.30em}

\subsection{What this costs you}

The sponsor pays for the drug, the platform, study-only imaging and laboratory
work, protocol-required visits, research-related injury care, and travel to a
stated cap. The participant's usual arrangement covers standard surgical care
and the inpatient stay, exactly as it would outside the study.

Two items are not settled by this protocol. Lost income is not compensated. And
continued access to daraxonrasib after study closure is not guaranteed. Both are
stated in the table, both are marked, and neither is buried. The National Cancer
Institute's guidance on who pays for clinical trials identifies exactly these as
the questions participants most often fail to ask in advance
\cite{NCICosts2024}.

\vspace{0.30em}

\begin{xltabular}{\textwidth}{L{5.0cm} L{3.0cm} Y}
\toprule
\textbf{Item} & \textbf{Who pays} & \textbf{Note} \\
\midrule
\endfirsthead
\multicolumn{3}{@{}l}{\footnotesize\itshape Table continued from the previous page.}\\
\toprule
\textbf{Item} & \textbf{Who pays} & \textbf{Note} \\
\midrule
\endhead
\midrule
\multicolumn{3}{r@{}}{\footnotesize\itshape Table continued on the next page.}\\
\endfoot
\bottomrule
\endlastfoot
Daraxonrasib & Sponsor & Supplied for the whole study period \\
Robotic platform use & Sponsor & No charge reaches you at any point \\
Study-only imaging and laboratory work & Sponsor & Anything the protocol adds beyond routine care \\
Protocol-required visits & Sponsor & The 18 visits of Figure 24 \\
Research-related injury care & Sponsor & Stated in the consent form, not only here \\
Travel and lodging & Sponsor, to a cap & The cap is stated in your consent form \\
Standard surgical care & Your usual arrangement & Exactly as it would be outside the study \\
Standard inpatient stay & Your usual arrangement & The 8 to 10 day stay is standard care \\
Long-term survivorship care & Your usual arrangement & On the routine pathway after handback \\
Lost income & \textbf{Not settled} & This protocol does not compensate it; H. R. 9510 v5 would address it \\
Continued drug access after closure & \textbf{Not settled} & Not guaranteed by this protocol \\
\end{xltabular}

\vspace{0.30em}

\subsection{After your last visit}

Care is handed back to the participant's own oncology team in writing, and a
full copy of the record goes to both the participant and that team. Survivorship
follow-up returns to the routine pathway. The device support obligation ends at
study closure, and the figure draws that ending rather than omitting it. The data
vault remains retrievable for fifteen years. Study results are sent to the
participant whether or not they completed, which is a commitment about the
sponsor's behaviour after the participant has stopped being useful to the study.

The longer-term proposal, a national platform that would carry post-trial
obligations across studies rather than leaving each sponsor to define its own,
is developed separately \cite{natlpaiplatf}. Figure 30 draws it as a dotted edge
rather than a solid one, because it is a proposal and not a provision.

\begin{pafloat}
\begin{pafig}[diagrams (python)-type: lifecycle across three time zones]
\dgnode{dgtile}{-7.4}{0}{\glyphpill}{Daraxonrasib, supplied by the sponsor}{z1}
\dgnode{dgtile}{-7.4}{-2.7}{\glyphgear}{Robotic platform use, no charge to you}{z2}
\dgnode{dgtile}{-4.6}{0}{\glyphflask}{18 protocol visits, study-only items paid}{z3}
\dgnodew{dgtiled}{-4.6}{-2.7}{\glyphshield}{Research-related injury care, sponsor pays}{z4}
\dgnode{dgtile}{0.4}{0}{\glyphuser}{Care handed back in writing to your own team}{y1}
\dgnode{dgtile}{0.4}{-2.7}{\glyphdoc}{Full record copy to you and your oncologist}{y2}
\dgnodeg{dgtilek}{3.2}{0}{\glyphgear}{Device support obligation ends at closure}{y3}
\dgnode{dgtile}{3.2}{-2.7}{\glyphclock}{Survivorship follow-up, routine pathway}{y4}
\dgnode{dgtile}{8.2}{0}{\glyphdb}{Your data, retrievable on request, 15 years}{x1}
\dgnode{dgtile}{8.2}{-2.7}{\glyphchart}{Study results sent to you either way}{x2}
\dgnodeg{dgtilek}{11.0}{0}{\glyphpill}{Continued drug access, not guaranteed}{x3}
\dgnode{dgtile}{11.0}{-2.7}{\glyphdoc}{Long-term safety registry, voluntary}{x4}
\dgnodew{dgtiled}{-5.8}{-6.6}{\glyphlink}{Sponsor pays: drug, platform, study-only work, injury care, travel to a cap}{p1}
\dgnodeg{dgtileg}{1.8}{-6.6}{\glyphhand}{Your usual arrangement: standard surgical care and the inpatient stay}{p2}
\dgnodeg{dgtilek}{9.4}{-6.6}{\glyphdoc}{Not settled by this protocol: lost income, and continued drug access}{p3}
\begin{scope}[on background layer]
\node[dgcluster,fit=(z1)(z1l)(z2)(z2l)(z3)(z3l)(z4)(z4l)]  (T1) {};
\node[dgcluster2,fit=(y1)(y1l)(y2)(y2l)(y3)(y3l)(y4)(y4l)] (T2) {};
\node[dgcluster2,fit=(x1)(x1l)(x2)(x2l)(x3)(x3l)(x4)(x4l)] (T3) {};
\node[dgcluster2,fit=(p1)(p1l)(p2)(p2l)(p3)(p3l)]          (T4) {};
\end{scope}
\node[dgctitle,anchor=south west]  at (T1.north west) {during the study, day $-28$ to month 24};
\node[dgctitle2,anchor=south west] at (T2.north west) {the 12 months after your last visit};
\node[dgctitle2,anchor=south west] at (T3.north west) {beyond, and what remains available};
\node[dgctitle2,anchor=south west] at (T4.north west) {who pays};
\draw[dgedge]  (z1) -- (z3);
\draw[dgedge]  (z2) -- (z4);
\draw[dgedgeb] (z4) to[out=0,in=180,looseness=0.85] node[font=\tiny\sffamily,above,pos=0.5]{continues if related} (y1);
\draw[dgedge]  (y1) -- (y2);
\draw[dgedge]  (y1) to[out=0,in=180,looseness=0.85] (y3);
\draw[dgedge]  (y2) to[out=0,in=180,looseness=0.85] (y4);
\draw[dgedgek] (z3) to[out=20,in=160,looseness=0.85] node[font=\tiny\sffamily,above,pos=0.5]{ends} (y3);
\draw[dgedge]  (y2) to[out=-20,in=200,looseness=0.85] (x1);
\draw[dgedge]  (y4) to[out=0,in=180,looseness=0.85] (x2);
\draw[dgedgek] (y3) to[out=20,in=160,looseness=0.85] (x3);
\draw[dgedge]  (x2) to[out=20,in=200,looseness=0.85] (x4);
\draw[dgedged] (p1) to[out=120,in=-90,looseness=0.8] (z2.south);
\draw[dgedged] (p1) to[out=60,in=-100,looseness=0.8] (z4.south);
\draw[dgedged] (p2) to[out=70,in=-100,looseness=0.8] (y4.south);
\draw[dgedgek] (p3) to[out=70,in=-100,looseness=0.8] (x3.south east);
\draw[dgedgeb,dotted,line width=0.9pt] (p3) to[out=30,in=-70,looseness=0.85]
  node[font=\tiny\sffamily,right,pos=0.5]{H. R. 9510 v5 would close both} (x4.south east);
% the terminator bars that mark an obligation ending
\foreach \n in {y3,x3}{\draw[protoblack,line width=1.1pt] ([xshift=-1.5mm,yshift=-6.4mm]\n.south) -- ([xshift=1.5mm,yshift=-6.4mm]\n.south);}
\node[font=\scriptsize\sffamily\itshape,text=protogray,anchor=north west,align=left]
  at (-9.4,-8.6) {An obligation that ends is drawn ending, with a terminator bar, rather than quietly left out. Two items are marked not settled, and\\the legislative proposal that would close them is drawn dotted, because it is a proposal and not a provision of this protocol.};
\end{pafig}
\vspace{-0.7cm}
\figcaption{Figure 30. Post-trial continuity across three time zones, during the study, the 12\\
months after the last visit, and beyond, with the payer attached to every item and the\\
two unsettled items drawn as unsettled rather than being quietly left out of the figure.}
\end{pafloat}

\subsection{If you take one page to the consent conversation}

Everything above is the argument. What follows is the page to carry, because a
participant reading forty pages the night before a consent conversation will
remember five things at most and should choose which five.

\textbf{Ask who is accountable for the failure class you care about most, and
ask for the matrix.} It is published, it has one accountable party per row, and
if the answer names a committee rather than a person, the answer is wrong.

\textbf{Ask for the two numbers measured in your own operating room.} The
emergency-stop latency, cross-arm and system-wide, is measured for your case
before induction. They are the only performance numbers in this whole paper
that will be about you specifically.

\textbf{Ask what happens if you say no today and yes in six weeks.} The answer
should be that eligibility is re-assessed at the moment a slot opens, that
nothing about declining now prejudices later enrollment, and that your standard
care does not change either way.

\textbf{Ask which of your concerns is answered by governance alone.} There are
five, they are listed in \S 3.8, and a site that claims all twenty-one are
answered by hard limits has not read the same protocol this paper read.

\textbf{Ask for your audit extract to be demonstrated, not described.} The
extract exists from the evening of the operation. A site that can show the
mechanism before you consent is a site that will produce it afterwards.

Those five questions take about ten minutes and they test the five properties
the rest of this paper spends forty pages establishing. If the answers are
satisfactory, the argument of this paper is that participation is a reasonable
choice. If they are not, the argument of this paper is that the participant is
entitled to say so and to be believed, and \S 12.2 is the enumeration of rights
that makes saying so safe.

\subsection{Closing argument}

The argument of this paper is not that this trial is safe because it is new. It
is the opposite. Novelty is the reason to demand specifics, and \S 3 collects
twenty-one specific demands that people actually make. What the paper claims is
that this protocol answers more of them, more precisely, than the alternative
does.

Consider what a participant knows after reading it. They know the force each
robotic arm may apply and the force all eight may apply together. They know how
long stopping takes, in milliseconds, and who may press stop. They know that no
message travels from the model to the arms, and they can see the diagram in
which the arrow is absent. They know the read list on every store that holds
their data and how many years each is kept. They know the name of the person
accountable for each of nine failure classes and which of them to call. They
know the two costs the protocol does not cover. And they know that a Phase 1
study of eighteen participants cannot tell them whether this improves survival.

A participant undergoing a standard Whipple outside a trial knows almost none of
that, not because anyone is concealing it, but because nobody assembled it. That
asymmetry is the argument. It is not an argument that the technology is
harmless; it is an argument that a protocol which can be interrogated this
closely is a better place to be than one which cannot.

The honest gaps remain where \S 3.7 put them. Cancer-control effectiveness is
answered by a comparator and not by a result. Accountability is answered by a
matrix and not by a settled liability doctrine. The human relationship is
answered by a notification rule and not by a guarantee. Cost is answered line by
line with two lines left open. Sections 10, 11, and 12 say what would have to
change for each, and the proposed legislation \cite{hr9510billv5} and the
clinician-trust framework \cite{clintrustpai} are the routes by which some of it
might. Until then, this paper states the gaps in the same typeface as the
answers, because a proponent who marks their own limits is the only kind worth
believing \cite{onpremwhippl}.
